\documentclass[11pt,a4paper]{article}
\usepackage{harmonikabau_preamble}

% == Document Metadata ====================
\newcommand{\hbdoknr}{0500}
\newcommand{\hbtitel}{Guide to Aesthetic Forensics}
\newcommand{\hbuntertitel}{of Accordion Casework}
\newcommand{\hbheadtext}{Doc.\,0500 --- Aesthetic Forensics: Accordion Casework}
\newcommand{\hbfoottext}{Accordion Construction}

\AtBeginDocument{%
  \fancyhead[L]{\small\hbheadtext}%
  \fancyhead[R]{\small\thepage}%
  \fancyfoot[C]{\small\color{hb@darkgray}\hbfoottext}%
}

\begin{document}
\hbmaketitle

	\section*{Note on Methodology}
	This guide focuses exclusively on external design features (grille, veneer, fittings). Note that the technical classification (reed plates, mechanics, reed blocks) forms a separate chapter and must be considered independently of the outer casing. A historic casework may well contain later technical components.

	\section*{1. Art Nouveau Aesthetics (ca.\ 1895--1914)}
	Art Nouveau is not mere imitation of nature, but an abstraction of it. Its aim was to ennoble everyday objects through organic forms.

	\begin{itemize}
		\item \textbf{Line quality:} Dominance of the so-called ``whiplash line''. Forms appear fluid, undulating and asymmetric. Hard edges and mathematical grids are rare.
		\item \textbf{Motifs:} Floral inspiration: tendrils, blossoms and insects. Even when patterns are arranged symmetrically, they appear alive and ``grown''.
		\item \textbf{Grille character:} The grille functions as a decorative element that breaks the surface organically. Decoration is flat and ``breathing''.
	\end{itemize}

	\section*{2. Aesthetics of the 1920s (Art D\'eco)}
	After 1918, the spirit of the times changed fundamentally. Nature was replaced by the machine; design became urban and technocratic.

	\begin{itemize}
		\item \textbf{Geometric discipline:} Rectangles, triangles, stepped forms and the meander (Greek key) dominate. Geometry is not a filler but an architectural directive.
		\item \textbf{Industrial rhythm:} Patterns such as the meander or repeating ovals are designed for industrial mass production (stamping). They appear cool, predictable and precise.
		\item \textbf{Veneer design:} Walnut is used in large surfaces, not as a base for carvings. A simple black-and-white contrast border (stringing) underscores the geometry of the wood.
		\item \textbf{Fittings:} Use of cool nickel silver or nickel. Fittings now mark the corners rather than ``decorating'' them, which emphasises the construction of the instrument.
	\end{itemize}

	\section*{3. The Collector's Trap: Historicist Geometry}
	A common mistake is confusing pre-war geometry with Art D\'eco.

	\begin{itemize}
		\item \textbf{Pre-war period (geometric historicism):} Here the \textit{horror vacui} (fear of empty space) prevails. The surface is completely covered with small, often overloaded geometric inlays. Every element is individually fitted. This is a demonstration of craftsmanlike patience, not industrial efficiency.
		\item \textbf{Distinction:} Ask yourself: does the pattern look like a densely woven, hand-crafted carpet (pre-war) or like a clear, rhythmic skeleton (1920s)?
	\end{itemize}

	\section*{Conclusion}
	The decoration of an accordion is a mirror of its era. While Art Nouveau wanted nature to grow into the instrument, Art D\'eco sought to define the instrument as a technical, precise object of modernity. With floral motifs, focus not on the ``what'' (e.g.\ a lily) but on the ``how'' (is the line organically fluid or statically geometric?).


\end{document}
